\documentclass[11pt]{article}

\usepackage[T1]{fontenc}
\usepackage{lmodern}
\usepackage{microtype}
\usepackage[margin=1.02in]{geometry}
\usepackage{amsmath,amssymb,amsthm,mathtools}
\usepackage{bm}
\usepackage{booktabs}
\usepackage{authblk}
\usepackage{enumitem}
\usepackage[hidelinks]{hyperref}
\usepackage[nameinlink,noabbrev]{cleveref}

\usepackage[utf8]{inputenc}

\usepackage{tikz}
\usetikzlibrary{arrows.meta,calc,positioning,fit,backgrounds}

\usepackage{comment}
\usepackage{ulem}

\usepackage[noadjust]{cite}
\usepackage[
  convert-footnotes = true,
  placement = mixed
]{notes2bib}

\newtheorem{theorem}{Theorem}[section]
\newtheorem{lemma}[theorem]{Lemma}
\newtheorem{proposition}[theorem]{Proposition}
\newtheorem{corollary}[theorem]{Corollary}
\newtheorem{remark}[theorem]{Remark}
\newtheorem{definition}[theorem]{Definition}

\newcommand{\Cut}{\operatorname{c}}
\newcommand{\mc}{\operatorname{MC}}
\newcommand{\Eval}{\operatorname{Eval}}
\newcommand{\M}{\#\mc}
\newcommand{\FP}{\mathsf{FP}}
\newcommand{\ii}{\mathrm{i}}
\newcommand{\one}{\mathbf{1}}
\newcommand{\abs}[1]{\left\lvert #1\right\rvert}
\newcommand{\set}[1]{\left\{#1\right\}}

\newcommand{\sh}[1]{{#1}}

\title{Evaluating QAOA expectation values can be as hard as counting optimal solutions}
\author{Stuart Hadfield%
\thanks{\href{mailto:stuart.hadfield@nasa.gov}{stuart.hadfield@nasa.gov} \\ \sh{This document was created with assistance from  ChatGPT-5.6-Sol AI tools as part of the drafting of Figure~\ref{fig:triangle-gadget} and for verification of the proofs of this manuscript. All content has been reviewed and edited by the author.} }}
\affil{Quantum Artificial Intelligence Laboratory (QuAIL), NASA Ames Research Center, Moffett Field, CA 94035, USA}
\affil{USRA Research Institute for Advanced Computer Science, Mountain View, CA 94043, USA}

\date{\today}

\begin{document}
\maketitle

\begin{abstract}

Evaluating expectation values is a critical task for variational quantum eigensolvers, and for parameterized quantum circuits 
and other quantum algorithms more generally. 
We consider the well-studied case of the Quantum Approximate Optimization Algorithm (QAOA) for the MaxCut problem.
Recent work of Wang et al. [arXiv:2511.20212] showed this task to be NP-hard 
in general %
for any QAOA %
depth $p\geq 2$, %
complementing past results showing %
efficiently computable formulas for %
$p=1$ with arbitrary problem graphs. 
We sharpen this dichotomy showing that for
$p\geq2$ exact or exponentially precise 
cost expectation value 
evaluation is $\#$P-hard under deterministic polynomial-time Turing
reductions. %
Hardness %
at $p\geq 2$ is shown to remain even for evaluating 
single %
pairwise correlators %
$\langle Z\otimes Z\rangle$, 
as well as for
highly 
restricted sets of algorithm parameters. 
Our proof refines the NP-hardness construction of Wang et al. that recovers the maximum cut value from the largest exponent of a QAOA Laurent polynomial, utilizing a %
distinct %
and simpler construction that %
extracts a value proportional to the total number of maximum cuts, in addition to the optimal cut value. 
Thus we show that 
the QAOA expectation value hardness transition from
$p=1$ to $p=2$ is not only from tractability to %
optimization %
hardness, but to that of counting optimal solutions. As an application we show our results imply %
analogous hardness results for computing gradients and Hessians of QAOA circuits. 
\end{abstract}



\section{Introduction}


Quantum computing has %
attracted much interest toward enabling %
novel approaches to NP-hard combinatorial optimization problems, in both exact and approximate optimization settings~\cite{abbas2023quantum}. 
At the same time, rigorous results characterizing their power or limitations remain relatively sparse, %
impeding our assessment of their true 
utility or potential for quantum advantage here. %
A popular paradigm is the 
Quantum Approximate Optimization Algorithm~(QAOA)~\cite{Farhi2014,hadfield2019quantum}, in particular applied to the prototypical NP-hard MaxCut problem. 
For analyzing quantum algorithms such as QAOA it is common to bound its expected approximation ratio in terms of the expectation value of a quantum observable encoding the cost function. This corresponds to the expected cost value of samples obtained from repeated preparation and measurement of a fixed QAOA state. 
For QAOA at its lowest depth (number of layers) $p=1$ for MaxCut, exact analytic formulas for the expected cost value were obtained in \cite{Wang2018}. On the other hand, recent work of \cite{WangEtAl2025} has shown that computing the same quantities for $p\geq 2$ is 
itself in general NP-hard 
(i.e., can be as hard as the underlying problem itself). Our main contribution is to sharpen this dichotomy to show that the $p\geq 2$ case is in fact $\#$P-hard; i.e., the latter case is as hard as counting the number of optimal solutions of a constraint satisfaction problem, not just finding one, which is generally believed to be a harder %
task~\cite{Valiant1979,Krentel1988}. 




\paragraph{QAOA for MaxCut}
Consider the MaxCut problem on a graph $G=(V,E)$ on $n=|V|$ vertices.   
We identify cuts as spin (bit) assignments $x\in\{\pm1\}^{n}$ %
with cut value 
\begin{equation}
 \Cut_G(x)=\sum_{\{u,v\}\in E}\frac{1-x_ux_v}{2},
 \quad
 \mc(G)=\max_x\Cut_G(x),
 \quad
 \M(G)=\set{x:\Cut_G(x)=\mc(G)}.
 \label{eq:cut-definitions}
\end{equation} 
Finding a maximum cut %
is a canonical NP-hard optimization problem.  
Counting maximum cuts exactly, i.e. determining $|\M(G)|$, is $\#$P-hard~\cite[Lemma~13]{JerrumSinclair1993} (see also~\cite{provan1983complexity,hermann2009complexity}). 



For QAOA applied to MaxCut we define the standard cost and mixer Hamiltonians from the \sh{single-qubit} Pauli matrices $X=(\begin{smallmatrix} 0 & 1 \\ 1 & 0 \end{smallmatrix} )$ and
 $Z=(\begin{smallmatrix} 1 & 0 \\ 0 & -1 \end{smallmatrix} )$
as 
\begin{equation}
 C_G=\sum_{\{u,v\}\in E}\frac{I-Z_uZ_v}{2},
 \qquad
 B=\sum_{v\in V}X_v.
 \label{eq:hamiltonians}
\end{equation}
\sh{Here $C_G$ satisfies $C_G\lvert x\rangle=c_G(x)\lvert x\rangle$ for every computational basis state $\lvert x\rangle$. } %
At depth %
parameter~$p$, standard QAOA for MaxCut prepares the quantum state\footnote{\sh{Throughout, the QAOA depth parameter $p$ denotes the total number of
alternating cost and mixer rounds, rather than the physical gate
depth of an implementation.  Standard implementations of
$e^{-i\gamma C_G}$ %
require at least $\Delta(G)$ two-qubit %
gate depth. The hard instances we construct below 
have $\Delta(G)=N-1$, so their physical
two-qubit gate depth grows with the instance size,  even when
$p=2$.}} 
\begin{equation}
\lvert\psi_p(\boldsymbol\gamma,\boldsymbol\beta)\rangle
  =(e^{-\ii\beta_p B}e^{-\ii\gamma_p C_G})\dots (e^{-\ii\beta_2 B}e^{-\ii\gamma_2 C_G})( e^{-\ii\beta_1 B}e^{-\ii\gamma_1 C_G})
 \lvert+\rangle^{\otimes n},
 \label{eq:qaoa-state}
\end{equation}
\sh{for easy-to-prepare initial product state $\lvert+\rangle^{\otimes n}=(\lvert0\rangle+\lvert1\rangle)^{\otimes n}/\sqrt{2^n}$, }
while seeking $2p$ parameters (angles) 
$\gamma_1,\dots,\gamma_p,\beta_1,\dots,\beta_p$ optimizing the cost expectation value 
\begin{equation}
 F_p(G;\boldsymbol\gamma,\boldsymbol\beta)
 =\langle\psi_p\rvert C_G\lvert\psi_p\rangle.
 \label{eq:qaoa-objective}
\end{equation}
Given $F_p(G;\boldsymbol\gamma,\boldsymbol\beta)$ one can then obtain bounds on the expected approximation ratio obtained by %
repeatedly preparing and measuring the QAOA state 
with the same parameters on a quantum computer~\cite{Farhi2014,Wang2018,hadfield2018quantum}. \sh{Evaluations of $F_p(G;\boldsymbol\gamma,\boldsymbol\beta)$ 
are also important for other algorithm aspects including 
parameter optimization.}

The complexity of evaluating \eqref{eq:qaoa-objective} is distinct from the complexity of
sampling from the full QAOA output distribution~\cite{farhi2016quantum}.  At $p=1$, the expectation value has an explicit
formula on general graphs~\cite{Farhi2014,Wang2018} (see Appendix~\ref{sec:p1}.).  By contrast, Wang, Chen, Sun,
and Zhang recently showed~\cite{WangEtAl2025} that exact evaluation is NP-hard for every fixed $p\geq2$. %
Their hardness proof writes the cost expectation as a Laurent polynomial in a %
phase
variable and recovers $\mc(G)$ from its largest nonzero exponent when evaluated for $C_{G'}$ on a suitably transformed graph $G'=G'(G)$.


The present work upgrades this mechanism for $p\geq 2$ from optimization to counting hardness.  
We construct instances for
which %
the 
extremal %
Laurent polynomial coefficient can be evaluated in closed form and is shown to be proportional to the full
maximum-cut degeneracy $|\M(G)|$ of the input problem graph.  
The required coefficient uniformity is enforced by %
complete bipartite subgraph  
$K_{L,L}$ variable gadgets with suitably chosen $L=O(n)$,
along with densely connected anchor variables. 
See Figure~\ref{fig:triangle-gadget} below for an illustrative example of our construction. 
Our main theorem uses polynomially many expectation value queries at different phase values to interpolate the required Laurent coefficient for $p=2$.  
Hence we show Turing reducibility between these problems, written informally as 
\begin{eqnarray*}
    |\M| \; \leq_T\; \Eval^{\mathrm{MC}}_{p=2}.
\end{eqnarray*}
The general $p$ result then follows setting all angles in the additional $p-2$ layers to zero. 
This yields $\#$P-hardness for computing QAOA expectation values rather than NP-hardness as in prior work. %
We further show that hardness remains when restricted to computing a single-edge
correlator $\langle Z_rZ_s\rangle$. %






\subsection{Main result: Evaluating QAOA $p=2$ expectation values is \#P-hard}







\begin{definition}[QAOA expectation value evaluation problem]
For fixed \(p \in \mathbb{N}\), let $\Eval_p^{\mathrm{MC}}$ be the function problem that maps a  graph $G$ and %
parameters $\boldsymbol\gamma,\boldsymbol\beta$ to the exact algebraic number $$F_p(G;\boldsymbol\gamma,\boldsymbol\beta).$$


The additive-error version %
additionally receives %
$b\in\mathbb{N}$ in %
\sh{unary}\footnote{Here %
unary encoding of $b$ ensures that the requested number of accuracy bits
is part of the input length.  In particular, choosing \(b=O(N)\) permits
exponentially small error while preserving a polynomial-time reduction.}
and returns a rational number
$\widetilde F$ satisfying
$$\bigl|\widetilde F-
F_p(G;\boldsymbol{\gamma},\boldsymbol{\beta})\bigr|
\leq 2^{-b}.$$ 
\end{definition}

For a designated edge \(e=\{u,v\}\in E(G)\), we additionally define
\[
  \mathrm{Eval}^{\mathrm{uv}}_p
  (G;\boldsymbol{\gamma},\boldsymbol{\beta})
  :=
  \langle\psi_p|Z_uZ_v|\psi_p\rangle .
\]








We will primarily be concerned with deterministic polynomial-time Turing reductions, i.e., what we can solve efficiently given a black box computing $\Eval_p^{\mathrm{MC}}$. 
The output of an exact evaluation oracle is understood in a standard basis of the compositum
of the input number fields.  The particular representation is inessential for the Turing
reductions below because every query uses a number field of polynomial degree and the inverse
Fourier transform is carried out exactly within that field.
QAOA angles are specified exactly by algebraic values of their sines and cosines; the reduction
uses only roots of unity and fixed quadratic algebraic numbers.  Hence the oracle output is %
assumed to be a standard exact representation of elements of the %
 generated number field.\footnote{\sh{ %
For completeness, the oracle calls used in the reduction of Thm.~\ref{thm:main} may be
formalized as follows.  Set $Q=4m+1$ and
$\zeta_Q=e^{2\pi i/Q}$.  The phase $\phi_j=2\pi j/Q$ is encoded by
the integer pair $(Q,j)$, with the convention
$e^{i\phi_j}=\zeta_Q^j$, while the fixed mixer angles are encoded by
$(\cos\beta_1,\sin\beta_1)
   =(\tfrac2{\sqrt{5}},\tfrac1{\sqrt{5}})$ and $(\cos\beta_2,\sin\beta_2)
   =(\tfrac{1}{\sqrt{2}},\tfrac{1}{\sqrt{2}})$.
The oracle output is represented as an element of
$\mathbb{K}_Q
   =\mathbb{Q}(\zeta_Q,i,\sqrt{2},\sqrt{5})$
under the canonical embedding
$\zeta_Q\mapsto e^{2\pi i/Q}$, $i\mapsto\sqrt{-1}$, and with the
positive choices of $\sqrt{2}$ and $\sqrt{5}$.  Since
$[\mathbb{K}_Q:\mathbb{Q}]\leq 8\varphi(Q)$ and $Q$ is polynomial
in the instance size, exact arithmetic, equality testing, and the
inverse discrete Fourier transform 
in the proof of Thm.~\ref{thm:main} 
can all be performed in polynomial time
and with polynomially scaling number of bits~\cite{lenstra1982factoring,basu2006algorithms,cohen2013course}.  Any standard exact
number-field representation supporting these operations is
sufficient.} 
}




\begin{theorem}%
[QAOA expectation complexity dichotomy]\label{thm:main}
For unweighted simple graphs,
\begin{eqnarray}
 &&\Eval^{\mathrm{MC}}_1\in\FP,\\
 &&\Eval^{\mathrm{MC}}_p\text{ is }\#\text{P-hard for every fixed }p\geq2.
 \label{eq:threshold}
\end{eqnarray}
Moreover, for every fixed $p\geq2$,
it suffices to restrict evaluations to QAOA angles
\begin{equation}
 \gamma_1=\gamma_2, %
 \qquad
 \beta_1=\arctan \frac12,
 \qquad 
 \beta_2=\frac{\pi}{4},
 \qquad
 \gamma_j=\beta_j=0\quad(3\leq j\leq p).
 \label{eq:base-restrictions}
\end{equation}
\end{theorem}
\begin{corollary}
    The same conclusions of the theorem hold if the oracle %
evaluates only one designated two-qubit correlator
$\langle Z_rZ_s\rangle$ rather than the global cost expectation value $\langle C_G \rangle $.
\end{corollary}



Here the mixer angles are fixed constants independent of the input, and 
only the single tied cost phase angle 
$\phi:=\gamma_1=\gamma_2$ 
is varied between oracle queries. 



We show how the hardness dichotomy of the theorem extends directly to the setting of exponentially small additive error in Section~\ref{sec:expoApprox}.












\begin{remark}[Consequences of efficient exact evaluation] 
 Suppose that, for %
 $p\ge 2$, the exact \sh{or exponentially precise} expectation value of
the MaxCut QAOA circuit can be computed in %
polynomial
time on arbitrary graphs. %
Then by Toda's theorem~\cite{Toda1991} the polynomial hierarchy collapses to its lowest level, i.e. P=PH.
\end{remark}



For exponentially precise evaluation, we show below that, for
an $N$-vertex oracle instance, the extreme Laurent coefficient
has magnitude at least %
$2^{-cN}$ for a known constant $c$, and hence $O(N)$ bits of
additive precision suffice to recover the exact integer count
by rounding.


We emphasize that both our results and those of Ref.~\cite{WangEtAl2025} imply that there exist \sh{worst-case} families of graphs for which evaluating $F_p$ appears computationally intractable, but this by no means implies this property applies to all graphs. %
Indeed, Ref.~\cite{WangEtAl2025} establishes complementary 
algorithms parameterized by local treewidth~\cite{WangEtAl2025}, complementing 
results obtained for other special cases %
for which %
rigorous results have been obtained
in the literature \cite{Farhi2014,Farhi2014b,Wang2018,hadfield2018quantum,jiang2017near,bravyi2020obstacles,hastings2019classical,marwaha2021local,wurtz2021maxcut,barak2021classical,marwaha2022bounds,hadfield2022analytical,ozaeta2022expectation,basso2021quantum,boulebnane2021predicting,farhi2022quantum,basso2022performance,boulebnane2025evidence}. 
Furthermore, whether one can efficiently estimate a given QAOA cost expectation value on a classical computer or not, a quantum computer %
may still be required to obtain corresponding problem solution samples~\cite{farhi2016quantum}. 






\subsection{Relationship to the results of \cite{WangEtAl2025} and related work}

Prior to \cite{WangEtAl2025}, a sequence of papers \cite{Farhi2014,Farhi2014b,Wang2018,hadfield2018quantum,jiang2017near,bravyi2020obstacles,hastings2019classical,marwaha2021local,wurtz2021maxcut,marwaha2022bounds,ozaeta2022expectation,hadfield2022analytical,basso2021quantum,boulebnane2021predicting,farhi2022quantum,basso2022performance,boulebnane2025evidence} obtained formulas or bounds for $\Eval_p^{\mathrm{MC}}$ in special cases. 
Beyond performance analysis, evaluating QAOA expectation values is also critical for algorithm parameter selection; in \cite{bittel2021training} it was shown that the %
parameter optimization problem in some settings can itself become NP-hard. 


For computing QAOA expectation values, Reference~\cite{WangEtAl2025} previously established
NP-hardness for every fixed $p\geq2$ by identifying the maximum cut value with the largest nonzero exponent
of a Laurent polynomial.  
Here we retain the Laurent %
polynomial approach
but replace the
underlying graph construction with one that utilizes a balanced counting gadget, and prove a corresponding exact leading  coefficient identity.  
This results in a simpler and more intuitive %
reduction than that of~\cite{WangEtAl2025}, utilizing fewer types of gadgets in our construction. 
The resulting new conclusions are: (i) $\#$P-hardness;
(ii) recovery of the number, rather than merely the value, of maximum cuts; (iii) hardness
for a single edge correlator; and (iv) the fixed-%
angles restrictions of Theorem~\ref{thm:main}. Our analysis is primarily concerned with the $p=2$ case. The %
$p\geq 2$ statement then follows by padding a
hard depth-two circuit with $p-2$ identity layers (i.e., layers for which $\gamma=\beta=0$).



Our results apply %
instead to computation of QAOA expectation values; 
these quantities 
can also be estimated \textit{quantumly} %
by repeated sampling from the quantum computer. 
In \cite{farhi2016quantum}
Farhi and Harrow %
adapted the postselection-based %
framework of \cite{bremner2011classical}
to show that efficient classical sampling from the output of standard %
QAOA circuits 
even at depth $p=1$ 
would imply a collapse of the polynomial hierarchy (to at least its third level), 
with this result further sharpened in the recent paper~\cite{abolicnvs2026sharp}. 
    Krovi later proved average-case hardness results~\cite{krovi2022average} for approximating output probabilities of random
    quantum circuits, including random $p=1$ QAOA ensembles.
In~\cite{dalzell2020many}
    the authors refine asymptotic supremacy arguments into a fine-grained framework, allowing 
    quantitative statements about the scale required for classical intractability based on standard conjectures in complexity theory. 
Hence in terms of sampling tasks, QAOA circuits provide a potential path to quantum computational supremacy~\cite{farhi2016quantum,harrow2017quantum,dalzell2020many}. Our complementary results show a stronger 
\sh{conditional} collapse of the polynomial hierarchy to PH=P if one could 
efficiently %
evaluate QAOA expectation values (to exponential precision), but with distinct threshold $p=2$. 






Our result should also be distinguished from work on more general
quantum mean value problems for constant-depth circuits
\cite{bravyi2021classical}.  
\sh{In that setting, depth refers to physical circuit depth in a local
gate model, and substantially broader circuit and observable
families are considered. 
By contrast the fixed parameter $p$ in
our theorems counts QAOA rounds.}
\sh{As the graphs produced by our
reduction contain universal vertices, a standard two-qubit-gate
implementation of a cost layer has depth $\Omega(N)$; consequently,
our theorem is not a hardness result for constant physical-depth
circuits.}
\sh{Its novelty is instead that exact $\#\mathrm{P}$-hardness
already occurs within the rigid family of standard fixed-round
MaxCut-QAOA circuits, with either the MaxCut Hamiltonian or one
designated two-qubit correlator as the observable.  Moreover, our
proof is qualitatively different, deriving hardness directly from
an extreme Laurent-polynomial coefficient rather than from
postselection-induced universality~\cite{bremner2011classical,bravyi2021classical}.  On the other hand, our theorem
concerns exact, or exponentially precise, evaluation and therefore
does not establish hardness in the constant-additive-error regime
considered in Ref.~\cite{bravyi2021classical}.}







Further afield, Szegedy~\cite{szegedy2019qaoa} previously proposed viewing QAOA %
expectation values as graph invariants, and investigated whether structural information 
such as graph non-isomorphism can be recovered from %
comparing these quantities. 
\sh{Distinct from the expectation value setting, in \cite{hadfield2021representation} it was show that 
$\#\mathrm{P}$-hardness can arise in the construction of
Hamiltonian representations for compactly specified Boolean
functions.} %
QAOA has also been proposed as a solution sampler
for approximate counting
\cite{DrapeauBanerjeeKourtis2025}.  
Finally, we note that beyond \cite{WangEtAl2025}, Laurent polynomials have also proven useful in two recent works solving QAOA for MaxCut on the special case of \lq\lq ring of disagrees\rq\rq\ (two-regular %
graphs)\cite{marwaha2026qaoa,kol2026machine}, confirming %
optimal expected cost scaling with $p$ originally conjectured in \cite{Farhi2014}. 



The following section gives the details and proof of Theorem~\ref{thm:main}. %
We show a direct application of our results to assessing the hardness of computing gradients and Hessians for QAOA circuits in Sec.~\ref{sec:derivative-hardness}, and provide some additional discussion of our results in Section~\ref{sec:discussion}.















\section{A balanced maximum-cut counting gadget}\label{sec:gadget}


\begin{figure}[h]
\centering
\includegraphics[width=0.85\textwidth]{fig_v2.png}
\caption{The balanced counting gadget for the source instance $K_3$ (triangle graph).  %
Blue and red indicate the two labels of one source maximum cut and its corresponding lift.  Every $K_{L,L}$ variable
gadget is cut completely, so each $A_v$ and $B_v$ lie on opposite sides; the
synchronous matchings among the $A$-parts reproduce the three source edges.
The separated anchors $r$ and $s$ are adjacent to one another and each is
joined to every vertex in $W$.  Each %
part $A_v,B_v$
in panel~(b) %
contains $L$ vertices; only representative vertices and edges are shown. 
For $H=K_3$, the %
constructed graph~$G(H)$ has $L=9$, $N=56$, and %
maximum cut counts $|\M(G)|=2|\M(H)|=12$.
\sh{This figure was created with assistance from  ChatGPT-5.6-Sol AI tools.}}
\label{fig:triangle-gadget}
\end{figure}

Let input graph $H=(V,E)$ have $n$ vertices and maximum degree $\Delta$.  
Define the odd integer
\begin{equation}
 L=2(\Delta+2)+1.
 \label{eq:L}
\end{equation}
For each $v\in V( H)$, introduce disjoint vertex sets $A_v=\{a_{v,1},…,a_{v,L}\}$ and $B_v$, each of size $L$, and add all
edges of the complete bipartite graph between them.  For each
$\{u,w\}\in E(H)$ and each $j\in[L]$, add the synchronous edge
$\{a_{u,j},a_{w,j}\}$ connecting %
the corresponding vertices in $A_u$ and $A_w$.  (Synchronous edges are not added for vertices in the $B_v$.) Finally add two anchor vertices $r,s$, the edge $e_\star=\{r,s\}$, and
join each anchor to every vertex in
\begin{equation}
 W=\bigcup_{v\in V(H)}(A_v\cup B_v).
\end{equation}
Denote the resulting graph by $G=\mathcal G( H)$.  It is unweighted, simple, and
connected, with
\begin{equation}
 N:=|V(G)|=2 nL+2 
 \label{eq:N}
\end{equation}
\sh{vertices. The number of edges is
\begin{equation}
  |E(G)|
  =
  nL^2+L|E(H)|+4nL+1.
\end{equation}
Since \(\Delta\le n-1\), we have \(L=O(n)\), and hence
 $N=|V(G)|=O(n^2)$ and 
 $ |E(G)|=O(n^3)$. 
Thus \(G(H)\), and all oracle queries generated from it, can be
constructed in polynomial time.  }

\begin{lemma}[Gadget rigidity]\label{lem:rigidity}
Every maximum cut of $G$ places each $A_v$ and $B_v$ monochromatically on opposite sides,
places $r$ and $s$ on opposite sides, and induces a maximum cut of $H$ through the
orientation of the sets $A_v$.  Conversely, every maximum cut of $H$ and either
orientation of $(r,s)$ determine a maximum cut of $G$.  Consequently,
\begin{align}
 \mc(G)&= nL^2+2 nL+1+L\mc( H),
 \label{eq:MCG}\\
 |\M(G)|&=2|\M(H)|.
 \label{eq:MG}
\end{align}
Moreover every maximum cut contains exactly $|W|/2= nL$ vertices of $W$ on each
side.
\end{lemma}

\begin{proof}
Each vertex of a variable $K_{L,L}$ has at most $\Delta+2$ external neighbors.  
Consider one variable gadget with parts $A_v$ and $B_v$.  Relative
to the nearer of its two proper bipartite orientations, suppose that
$a$ vertices of $A_v$ and $b$ vertices of $B_v$ are misplaced, and set
$q=a+b$.  %
The Hamming distances from the current labeling of
\(A_v\cup B_v\) to the two proper bipartite orientations sum to \(2L\).
Therefore the nearer orientation has Hamming distance
\(q\le L\).
The
number of cut edges internal to the gadget is
\[
(L-a)(L-b)+ab=L^2-Lq+2ab.
\]
Thus correcting the orientation of the gadget gains at least 
\[
Lq-2ab
   \geq Lq-\frac{q^2}{2}
   \geq \frac{Lq}{2}
\]
internal cut edges, where we used $2ab\leq q^2/2$ and $q\leq L$.
Each moved vertex has at most $\Delta+2$ external neighbors, so this
correction loses at most $(\Delta+2)q$ external cut edges.  
\sh{This bound counts incidences between moved vertices and external edges;
an external edge incident to two moved vertices may therefore be counted
twice, but this is harmless for an upper bound.} 
Since
$L>2(\Delta+2)$, the net change is strictly positive whenever $q>0$.
Consequently every maximum cut orients each $K_{L,L}$ gadget
properly.
Thus each gadget
contributes $L^2$ cut edges and exactly $L$ vertices to each side.

If the anchors lie on the same side, flipping one anchor leaves the total number of cut
anchor--$W$ edges unchanged, because $W$ is balanced, and additionally cuts $\{r,s\}$.
Hence the anchors are opposite.  With opposite anchors, every $w\in W$ contributes exactly
one cut anchor edge.  The only orientation-dependent edges are the synchronous edges:
the $L$ edges representing $\{u,v\}\in E( H)$ are cut exactly when the orientations
of $A_u$ and $A_v$ differ.  Equations \eqref{eq:MCG} and \eqref{eq:MG} follow immediately.
\end{proof}


The Hamming distances from a labeling of \(A_v\cup B_v\) to the two
proper bipartite orientations sum to \(2L\); hence the nearer
orientation has distance \(q\le L\).  Correcting all \(q\) misplaced
vertices simultaneously changes only external edges incident to those
vertices.  Bounding the possible loss by \((\Delta+2)q\) counts
incidences rather than distinct edges and may therefore overcount,
which is harmless for an upper bound.

Suppose a cut does not respect one of the balanced complete bipartite
gadgets \(A_v\cup B_v\). The Hamming distances from the induced labeling
to the two proper bipartite orientations sum to \(2L\); hence one of
these orientations differs on at most \(L\) vertices. Let \(q\le L\)
denote this minimum Hamming distance. Correcting those \(q\) vertices
restores all \(L^2\) gadget edges while changing only external edges
incident to the corrected vertices. The possible loss is at most
\((\Delta+2)q\), ...


Our graph construction is simpler and more symmetric than the gadget used in Ref.~\cite{WangEtAl2025}.  The latter employs asymmetric complete-bipartite vertex gadgets, a separate large complete-bipartite frame, and an additional controller vertex in order to isolate the maximal Laurent exponent.  By contrast, our reduction uses only balanced ($K_{L,L}$) variable gadgets and two universal anchors (r,s). %





\subsection{The extreme coefficient at depth two}\label{sec:coefficient}

Fix an edge $e=\{u,v\}$ of a graph $G$. 
Set $\gamma_1=\gamma_2=\phi$, 
$\beta_1=\beta$, $\beta_2=\pi/4$, and define
\begin{equation}
  H_e(e^{\ii\phi},\beta)
 :=\langle\psi_2(\phi,\phi;\beta,\pi/4)|\,Z_uZ_v\,|
 \psi_2(\phi,\phi;\beta,\pi/4)\rangle.
 \label{eq:edge-correlator}
\end{equation}
For $a,b\in\{\pm1\}$ let
$ K_\beta^\pm(a,b)=\langle a|e^{\pm\ii\beta X}|b\rangle$. 
Writing $c=\cos\beta$ and $s=\sin\beta$ we have
\begin{equation}
 K_\beta^\pm(a,b)=
 \begin{cases}c,&a=b,\\ \pm\ii s,&a\ne b,\end{cases}
 \label{eq:kernel-values}
\end{equation}
Let $T_e$ operate on an assignment string so as to flip the two %
vertex spins in $e$.




\begin{lemma}[Reduced path sum]\label{lem:path-sum}
For $z=e^{\ii\phi}$,
\begin{align}
 H_e(z;\beta)
 ={}&-2^{-N}\sum_{x,y,t\in\{\pm1\}^{V(G)}}
 y_uy_v\,
 z^{\Cut_G(x)-\Cut_G(t)+\Cut_G(y)-\Cut_G(T_ey)}
 \notag\\
 &\hspace{1.2cm}\times
 \prod_{w\in V(G)}K_\beta^+(x_w,y_w)
 K_\beta^-((T_ey)_w,t_w).
 \label{eq:path-sum}
\end{align}
Hence $H_e(z;\beta)$ is a Laurent polynomial in $z$.
\end{lemma}

\begin{proof}
Insert computational-basis resolutions between all unitaries.  The two cost layers on the
ket and bra contribute the exponent in \eqref{eq:path-sum}.  Summing over the intermediate
spins adjacent to the final mixer gives, for $w\notin e$,
\begin{equation}
 \sum_{q_w}K_{\pi/4}^+(y_w,q_w)K_{\pi/4}^-(q_w,\bar y_w)
 =\delta_{y_w,\bar y_w},
\end{equation}
and, for $w\in e$,
\begin{equation}
 \sum_{q_w}q_wK_{\pi/4}^+(y_w,q_w)K_{\pi/4}^-(q_w,\bar y_w)
 =-\ii y_w\delta_{\bar y_w,-y_w}.
\end{equation}
The endpoint product is $-y_uy_v$ and enforces $\bar y=T_ey$, yielding
\eqref{eq:path-sum}.
\end{proof}

Define the maximal boundary increment
\begin{equation}
 R_e=\max_y\bigl[\Cut_G(y)-\Cut_G(T_ey)\bigr].
 \label{eq:boundary}
\end{equation}

\begin{lemma}[Boundary increment]\label{lem:boundary}
For every edge $e=\{u,v\}$ of a simple graph,
\begin{equation}
 R_e=d_u+d_v-2.
 \label{eq:boundary-value}
\end{equation}
\end{lemma}

\begin{proof}
Only edges incident on exactly one endpoint can change when both endpoints are flipped, so
$R_e\leq d_u+d_v-2$.  Equality is attained by assigning the same spin to $u,v$ and the
opposite spin to every other neighbor of either endpoint. This assignment is consistent even when \(u\) and \(v\) have common
neighbors, since every vertex in
\((N(u)\cup N(v))\setminus\{u,v\}\) receives the same spin opposite
to \(u\) and \(v\).
\end{proof}

For the anchor edge $e_\star=\{r,s\}$ of $G=\mathcal G(H)$, both endpoints are
universal and therefore
\begin{equation}
 R_{e_\star}=2N-4.
 \label{eq:Rstar}
\end{equation}
The only maximizing boundary assignments are
\begin{equation}
 (q_\sigma)_r=(q_\sigma)_s=\sigma,
 \qquad
 (q_\sigma)_w=-\sigma\quad(w\in W),
 \qquad \sigma\in\{\pm1\},
 \label{eq:q-sigma}
\end{equation}
for which $T_{e_\star}q_\sigma=-\sigma\one$.  Since every nonanchor vertex has degree
strictly less than $N-1$, $e_\star$ uniquely maximizes $R_e$ among all edges. 
Hence every exponent in \eqref{eq:path-sum} is at most
\begin{equation}
 D_\star=\mc(G)+2N-4.
 \label{eq:Dstar}
\end{equation} 
%
\sh{Using \(c_G(t)\ge0\) the exponent in Eq.~\eqref{eq:path-sum} decomposes as
\[
c_G(x)-c_G(t)
+\bigl(c_G(y)-c_G(T_{e_\star}y)\bigr)
\le
\mathrm{MC}(G)+R_{e_\star}
=
D_\star.
\]
Since each summand is bounded
independently, } %
attaining the upper bound requires $c_G(x)=\mathrm{MC}(G)$, $c_G(t)=0$, and $c_G(y)-c_G(T_{e_\star}y)=R_{e_\star}$.  Since
$G=G(H)$ is connected, its only zero-cut assignments are the two
constant assignments.  
Hence
at equality, necessarily
\begin{equation}
 x\in\M(G),
 \qquad
  t=\tau\one,\,\,\tau \in \{\pm 1\},
 \qquad
 y\in\{q_+,q_-\}.
 \label{eq:extreme-histories}
\end{equation}
By Lemma~\ref{lem:rigidity}, each $x\in\M(G)$ is balanced on $W$ and separates the anchors.
Consequently
\begin{equation}
 d_{Hamm}(x,q_\sigma)=1+\frac{|W|}{2}=\frac N2, %
 \label{eq:hamming}
\end{equation}
independently of $x$ and $\sigma$.  


\begin{theorem}[Extreme-coefficient identity]\label{thm:coefficient}
Recall $D_* := \mathrm{MC}(G)+2N-4$ as in Eq.~\eqref{eq:Dstar}. Write
\begin{eqnarray}
    H_{e_\star}(z;\beta)=\sum_{k\in\mathbb Z}a_k(\beta)z^k.
\end{eqnarray}
For \(G=G(H)\), we have
  $a_k(\beta)=0$ for every $k>D_*$, 
and the coefficient at \(D_\star\) is
\begin{equation}
 a_{D_*}(\beta)
 =-2^{1-N}|\M(G)|\,
 c^{N/2}(\ii s)^{N/2}\left[c^N+(-\ii s)^N\right],
 \label{eq:edge-top-coefficient}
\end{equation}
where \(c=\cos\beta\) and \(s=\sin\beta\). %
In particular, if \(\tan\beta=1/2\), then
\(a_{D_\star}(\beta)\neq0\), and so 
\[
  D_\star=\max\{k\in\mathbb Z:a_k(\beta)\neq0\}.
\]
Hence, for this mixer angle \(|\#\mathrm{MC}(G)|\) is recovered by division
by a known nonzero algebraic factor.
\end{theorem}


\begin{proof}
Only the paths in Eq.~\eqref{eq:path-sum} satisfying the equality conditions in
Eq.~\eqref{eq:extreme-histories} can contribute to~$z^{D_*}$. 
The exponent in \eqref{eq:path-sum} attains $D^\star = \mathrm{MC}(G) + 2N-4$
if and only if $c_G(x) = \mathrm{MC}(G)$, $c_G(t) = 0$, and
$c_G(y) -c_G(T_{e^\star} y) = 2N-4$ simultaneously. %

Fix such a path $(x, q_\sigma, \tau\mathbf{1})$ with
$x \in \#\mathrm{MC}(G)$ and $\sigma, \tau \in \{\pm 1\}$. 
First observe the endpoint factor is
$y_r y_s = (q_\sigma)_r (q_\sigma)_s = \sigma^2 = 1$ for both $\sigma$. 
Next, by \eqref{eq:hamming} the Hamming distance
$h = d_{\mathrm{Hamm}}(x, q_\sigma) = N/2$ is the same integer for every
$x \in \#\mathrm{MC}(G)$ and both $\sigma$, so by \eqref{eq:kernel-values} we have
\begin{equation}
\prod_{w \in V(G)} K^{+}_{\beta}\bigl(x_w, (q_\sigma)_w\bigr)
= c^{\,N-h}(\mathrm{i}s)^{h}
= c^{N/2}(\mathrm{i}s)^{N/2},
\label{eq:Kplus-product}
\end{equation}
independently of $x$ and $\sigma$.
Finally, by \eqref{eq:extreme-histories} we have %
$t_w = \tau$ for all $w$, and since $T_{e^\star} q_\sigma = -\sigma\mathbf{1}$,
\begin{equation}
\prod_{w \in V(G)} K^{-}_{\beta}\bigl((T_{e^\star} q_\sigma)_w,\, t_w\bigr)
= K^{-}_{\beta}(-\sigma,\tau)^{N}
=
\begin{cases}
c^{N} & \tau = -\sigma\\[2pt]
(-\mathrm{i}s)^{N} & \tau = \sigma.
\end{cases}
\label{eq:Kminus-product}
\end{equation}
Summing \eqref{eq:Kminus-product} over $\tau \in \{\pm 1\}$ gives
$c^{N} + (-\mathrm{i}s)^{N}$, independently of $\sigma$. 
Since the endpoint factor, the product \eqref{eq:Kplus-product}, and the
$\tau$-sum are all independent of $x$ and $\sigma$, the triple sum over the
surviving paths factorizes as
\begin{align}
a_{D_*}(\beta)
&= -2^{-N} \sum_{x \in \#\mathrm{MC}(G)} \; \sum_{\sigma \in \{\pm 1\}}
   \; \sum_{\tau \in \{\pm 1\}}
   y_r y_s \prod_{w} K^{+}_{\beta}\bigl(x_w,(q_\sigma)_w\bigr)
   \, K^{-}_{\beta}\bigl((T_{e^\star} q_\sigma)_w, t_w\bigr) \notag\\
&= -2^{-N} \cdot \underbrace{|\#\mathrm{MC}(G)|}_{x\text{-sum}}
   \cdot \underbrace{2}_{\sigma\text{-sum}}
   \cdot \underbrace{1}_{y_r y_s}
   \cdot \underbrace{c^{N/2}(\mathrm{i}s)^{N/2}}_{\eqref{eq:Kplus-product}}
   \cdot \underbrace{\bigl[c^{N} + (-\mathrm{i}s)^{N}\bigr]}_{\tau\text{-sum}}
   \notag, %
\label{eq:extreme-coeff-derivation}
\end{align}
which rearranges to give Eq.~\eqref{eq:edge-top-coefficient}. 
For $\tan\beta=1/2$, the two terms in the bracket always have unequal
magnitudes and cannot cancel, which implies $a_{D_*}(\beta)\neq 0.$
\end{proof}






For the global expectation,
\begin{equation}
 \mathcal{F}_2(G;z,\beta):=F_2(G;\phi,\phi,\beta,\pi/4)=\frac{|E(G)|}{2}-\frac12\sum_{e\in E(G)}H_e(z;\beta).
 \label{eq:global-sum}
\end{equation}
Since $e_\star$ uniquely maximizes $R_e$, the coefficient of $z^{D_\star}$ in
\eqref{eq:global-sum} is exactly $-a_{D_\star}/2$.  Thus either the single correlator or
the global objective contains the same count information.



\subsection{Coefficient recovery and proof of the main theorem}\label{sec:hardness}

Let $m=|E(G)|$.  Every Laurent exponent in \eqref{eq:path-sum} lies in $[-2m,2m]$.  Set
\begin{equation}
 Q=4m+1,
 \qquad
 \omega=e^{2\pi\ii/Q},
 \qquad
 \phi_j=\frac{2\pi j}{Q}
 \quad(0\leq j<Q).
 \label{eq:dft-points}
\end{equation}
For either $P(z)=H_{e_\star}(z;\beta)$ or $P(z)=F_2(G;z,\beta)$, exact oracle values at
$z=\omega^j$ recover every coefficient by the inverse discrete Fourier transform
\begin{equation}
 a_k(\beta)=[z^k]P(z)=\frac1Q\sum_{j=0}^{Q-1}P(\omega^j)\omega^{-jk}.
 \label{eq:idft}
\end{equation}
%
\sh{%
As the Laurent exponent interval \([-2m,2m]\) has width
\(4m<Q\), two distinct exponents in this interval are never congruent
modulo \(Q\).  Hence the inverse discrete Fourier transform %
does not suffer from aliasing and recovers the ordinary Laurent coefficients, including
those with negative indices.}

\begin{proof}[Proof of Theorem~\ref{thm:main}]
The $p=1$ statement is Proposition~\ref{prop:p1}.  For $p=2$, start from an instance $H$ of counting
maximum cuts. 
Set
\begin{equation}
 \cos\beta_1=\tfrac2{\sqrt5},
 \qquad
 \sin\beta_1=\tfrac1{\sqrt5},
\end{equation}
which satisfies $\tan \beta_1=1/2$. 
The coefficient \eqref{eq:edge-top-coefficient} is nonzero for every even $N$ and directly
determines $|\M(G)|=2|\M(H)|$.  
\sh{%
All coefficients are obtained from Eq.~\eqref{eq:idft}, and the largest exponent with nonzero coefficient is selected. Theorem \ref{thm:coefficient} guarantees that this exponent is $D_*$
and that its coefficient is nonzero.}
Thus exact evaluation is $\#$P-hard at QAOA depth two %
for either the global expectation or the designated correlator.

For any fixed $p>2$, set all angles in layers $3,\ldots,p$ to zero.  These layers are the
identity, so the depth-$p$ oracle value equals the depth-two value.  This proves the
$\#$P-hardness statement for every fixed $p\geq2$.
\end{proof}








\begin{remark}
    The reduction of Theorem~\ref{thm:main} recovers more than the count of maximum cuts: the extreme exponent $D_* = MC(G) + 2N - 4$ itself reveals $MC(G)$, and hence $MC(H)$ via Eq.~\eqref{eq:MCG}, before the coefficient is examined. Thus a single run of the interpolation recovers both the maximum cut value and the number of maximum cuts of H, so the present construction also reproduces the NP-hardness conclusions of \cite{WangEtAl2025} directly.
\end{remark}



\paragraph{Hardness on diameter-two graphs}

\begin{corollary}%
\label{cor:diameter-two-hardness}
For every fixed \(p \ge 2\), exact evaluation of the MaxCut QAOA
expectation is \(\#\mathrm{P}\)-hard under deterministic
polynomial-time Turing reductions, even when restricted to connected,
unweighted, simple graphs of diameter two. The same conclusion holds
for the expectation of the designated two-qubit observable \(Z_r Z_s\),
where \(\{r,s\}\) is an edge joining two universal vertices.
\end{corollary}





\begin{proof}%
The graph \(G(H)\) constructed in the proof of
Theorem~\ref{thm:main} is unweighted, simple, and connected. The two
anchor vertices \(r\) and \(s\) are adjacent to one another and to
every vertex in every variable gadget. 
Let \(u,v \in V(G(H))\). If neither vertex is \(r\), then both are
adjacent to \(r\), and hence there is a path \(u-r-v\) of length at
most two. If one of the vertices is \(r\), then the pair is adjacent. 
Thus every pair of distinct
vertices is joined by a path of length at most two, so
$ \operatorname{diam}(G(H)) \le 2$.
The graph is not complete: for example, two distinct vertices in the
same part \(A_v\) of a variable gadget \(K_{L,L}\) are nonadjacent.
Hence 
$\operatorname{diam}(G(H)) = 2$.

\end{proof}



\begin{remark}
The diameter-two restriction is obtained at the expense of large
degree that scales with the problem size\sh{, as it must: any diameter-two graph on $N$ vertices has
maximum degree at least $\sqrt{N-1}$} (because from any vertex, at most \(1+\Delta+\Delta(\Delta-1)
=\Delta^2+1\) vertices can lie within distance two).
Each anchor is adjacent to every gadget vertex, and the
vertices within each \(K_{L,L}\) gadget have degree growing with the
source-instance 
input problem size. Consequently, the corollary does not establish
hardness for bounded-degree graph families.
\end{remark}





\subsection{Hardness of exponentially precise additive approximation} \label{sec:expoApprox}
Recalling the setting of Theorem~\ref{thm:main} let 
\begin{eqnarray}
      F_G(\phi)
  :=
  \langle \psi_p(\phi) | C_G | \psi_p(\phi) \rangle,
\end{eqnarray}
where $|\psi_p(\phi)\rangle$ is the single-parameter depth-$p$ QAOA state with $\gamma_1=\gamma_2=\phi$, $\beta_1=\textrm{arctan}(1/2), \beta_2=\pi/4$, and all other angles fixed $\gamma_j=\beta_j=0$.


\begin{corollary}[Exponentially precise additive evaluation]\label{cor:precision}
Fix \(p\ge2\) and %
mixer angle choices of 
Theorem~\ref{thm:main}.  There exists a constant \(\alpha>0\) such that
the following problem is $\#$P-hard under deterministic
polynomial-time Turing reductions: given an \(N\)-vertex MaxCut
instance from the family constructed in the proof and an exactly
represented tied phase angle \(\phi\), output a number
\(\widetilde F_G(\phi)\) satisfying
\[
   \bigl|\widetilde F_G(\phi)-F_G(\phi)\bigr|
   \le 2^{-\alpha N}.
\]
The same statement holds for the designated correlator
\(\langle Z_rZ_s\rangle\).
\end{corollary}

\begin{proof}
Write
\[
   F_G(z)=\sum_{k=-D}^{D}a_kz^k,
   \qquad z=e^{i\phi}.
\]
Evaluate the approximation oracle at \(M>2D\) distinct \(M\)-th
roots of unity and apply the normalized inverse discrete Fourier
transform.  If every oracle value has additive error at most
\(\varepsilon\), then every recovered Laurent coefficient has
additive error at most \(\varepsilon\).  %
\sh{
By Eq.~\eqref{eq:global-sum} and Theorem~\ref{thm:coefficient}, %
for the coefficient of $z^{D^\star}$ in \(F_G(z)\) we may write}
\[
a_{D^\star} = -\tfrac12\kappa_N |\M(G)| = -\kappa_N |\M(H)|. 
\]






For the mixer angle choice $\tan\beta=1/2$, write
$c_\beta=\cos\beta=\frac{2}{\sqrt{5}}$
and 
$s_\beta=\sin\beta=\frac{1}{\sqrt{5}}$.
Then %
$|\kappa_N|$ is bounded
from below (cf. Eq.~\eqref{eq:edge-top-coefficient}) by
\begin{align}
2^{1-N}(c_\beta s_\beta)^{N/2}
 \left|c_\beta^N+(-is_\beta)^N\right|
\nonumber
\,\geq&\,\,
2^{1-N}\left(\tfrac{2}{5}\right)^{N/2}
 \left(c_\beta^N-s_\beta^N\right)
\nonumber\\
\,=&\,\,
2^{1-N}\left(\tfrac{2}{5}\right)^{N/2}
 \tfrac{2^N-1}{5^{N/2}}
\nonumber\\
\,\geq&\,\,
\left(\tfrac{\sqrt{2}}{5}\right)^N\,
=
2^{-N\log_2(5/\sqrt{2})}.
\end{align}
Here we used the %
triangle inequality and
$2^N-1\geq 2^{N-1}$.  
Hence $|\kappa_N|\ge 2^{-\alpha_0 N}$ for %
$\alpha_0:=\log_2(5/\sqrt2)$.  
%
Next, let $\alpha = \alpha_0+1$. 
If every oracle expectation value is returned with absolute error at most~$\varepsilon\leq 2^{-\alpha N}\leq |\kappa_N|/2$, then the normalized inverse Fourier transform recovers every
Laurent coefficient with complex-modulus error at most~$\varepsilon$.  
%
Let $\widehat a_k$ denote the recovered coefficients. For
$k>D_\star$ the true coefficient vanishes, so
$|\widehat a_k|\le\varepsilon<|\kappa_N|$. Using 
$|\#\mathrm{MC}(H)|\ge2$ %
(global spin reversal preserves the cut value) 
we have 
$$|\widehat a_{D_\star}|\ge|\kappa_N|\,|\#\mathrm{MC}(H)| -\varepsilon\ge\tfrac32|\kappa_N|>|\kappa_N|.$$ 
Hence $D_\star$ is
the largest exponent whose recovered coefficient exceeds
$|\kappa_N|$ in modulus, a threshold computable from $N$ and
$\beta$. 
Dividing $\widehat a_{D_\star}$ by the known
nonzero factor $-\kappa_N$ yields $|\#\mathrm{MC}(H)|$ to additive error
$\varepsilon/|\kappa_N|\le 2^{-N}<1/2$, so rounding to the
unique nearest integer recovers the counts $|\#\mathrm{MC}(G)|$ and $|\#\mathrm{MC}(H)|$ exactly.


\sh{%
The same arguments apply for the %
single correlator
\(\langle Z_rZ_s\rangle\). In that case %
Theorem~\ref{thm:coefficient} gives the coefficient directly without the %
\(-\tfrac12\) factor from Eq.~\eqref{eq:global-sum}. 
 Thus the same
threshold and rounding argument recovers $|\#\mathrm{MC}(G)|$ and $|\#\mathrm{MC}(H)|$ directly; only
the known multiplicative factor relating the extreme coefficient to
the count is changed.
}


\end{proof}







\section{Hardness of gradients and Hessian}
\label{sec:derivative-hardness}

Consider the setting of Theorem~\ref{thm:main} where we restrict to the
tied-phase line
\[
  F_G(\phi)
  :=
  \left\langle C_G\right\rangle_{\gamma_1=\gamma_2=\phi}.
\]
The same statements hold with $C_G$ replaced by the designated observable
$Z_rZ_s$.  As above,
\begin{equation}
  F_G(\phi)=\sum_{k=-D}^{D} a_ke^{ik\phi},
  \label{eq:compact-derivative-laurent}
\end{equation}
where $D\leq 2|E(G)|=\operatorname{poly}(|G|)$.  For the graph $G(H)$ produced by the
reduction %
we again use that in either case 
by Theorem~\ref{thm:coefficient} there is %
a $D_\star>0$ such that
  $a_{D_\star}(G)
  =
  \kappa(H)\,|\M(H)|$, with 
  $\kappa(H)\neq 0$ computable from \(N\) and the chosen
observable. \sh{The exponent \(D_\star\) itself need not be known in
advance, since it contains \(\mathrm{MC}(H)\).}

\begin{theorem}[Hardness of phase derivatives]
\label{thm:compact-derivative-hardness}
For every fixed $p\ge2$ and every fixed integer $r\ge1$, exact computation
of
\[
  F_G^{(r)}(\phi)=\frac{d^r}{d\phi^r}F_G(\phi)
\]
is $\#\mathrm{P}$-hard under deterministic polynomial-time Turing
reductions.  This holds for unweighted simple connected MaxCut instances 
for the parameter choices of  Theorem~\ref{thm:main}, and for both the
total MaxCut expectation and the correlator $\langle Z_rZ_s\rangle$.
\end{theorem}



To the best of our knowledge we are not aware of an
earlier \(\#\mathrm P\)-hardness result for exact derivatives of the
standard fixed-depth QAOA for MaxCut objective.

\begin{proof} %
Differentiating \eqref{eq:compact-derivative-laurent} gives
\[
  F_G^{(r)}(\phi)
  =
  \sum_{k=-D}^{D}(ik)^r a_ke^{ik\phi}.
\]
Choose $M>2D$, let $\omega=e^{2\pi i/M}$, and set
$\phi_j=2\pi j/M$.  Orthogonality of the $M$th roots of unity yields
\[
  \frac1M\sum_{j=0}^{M-1}
  F_G^{(r)}(\phi_j)\omega^{-jk}
  =
  (ik)^r a_k.
\]
Since $D_\star>0$ and \(a_{D_\star}\ne0\),  the derivative
coefficient $ (iD_\star)^r a_{D_\star}$ is nonzero. Consequently, the largest exponent with nonzero recovered
derivative coefficient is again \(D_\star\); no oracle for \(F_G\)
itself is required. 
The extreme coefficient is recovered exactly as
\[
  a_{D_\star}(G)
  =
  \frac{1}{(iD_\star)^rM}
  \sum_{j=0}^{M-1}
  F_G^{(r)}(\phi_j)\omega^{-jD_\star},
\]
where the value $D_*$ is recovered as in the proof of Thm.~\ref{thm:main}. 
Applying this identity to $G(H)$ and using
Theorem~\ref{thm:coefficient}
recovers
$|\M(H)|$ in polynomial time.
\end{proof}



\begin{corollary}[Hardness of the full gradient and Hessian]
\label{cor:compact-gradient-hessian}
Consider the depth-two QAOA for MaxCut objective
 $ \mathcal F_G(\gamma_1,\gamma_2,\beta_1,\beta_2)
  :=\langle C_G\rangle$. 
 Exact computation of the full gradient
$\nabla\mathcal F_G$ and of the full Hessian $\nabla^2\mathcal F_G$ is
$\#\mathrm{P}$-hard under deterministic polynomial-time Turing reductions.
The same conclusion holds for single correlators.
\end{corollary}

\begin{proof}[Proof of Corollary \ref{cor:compact-gradient-hessian}]
On the affine line
$\theta(\phi)=(\phi,\phi,\beta_1,\beta_2)$, the chain rule gives
\[
  F_G'(\phi)
  =
  v_\gamma^{\mathsf T}\nabla\mathcal F_G(\theta(\phi)),
  \qquad
  F_G''(\phi)
  =
  v_\gamma^{\mathsf T}\nabla^2\mathcal F_G(\theta(\phi))v_\gamma,
\]
for $v_\gamma=(1,1,0,0)^T$. 
The conclusion follows from
Theorem~\ref{thm:compact-derivative-hardness} with $r=1$ and $r=2$.
\end{proof}

\begin{remark}
The result establishes hardness of the tied-phase directional derivative
and curvature, and hence of returning the full gradient or Hessian.  It
does not by itself imply hardness of each individual gradient coordinate
or Hessian entry separately.
\end{remark}

\begin{corollary}[Exponentially precise derivative approximation]
\label{cor:approximate-derivative-hardness}
For %
the mixer angles of Theorem~\ref{thm:main} and every fixed $r\ge1$,
there exists a constant $c>0$ such that approximating
$F_G^{(r)}(\phi)$ to additive error at most $2^{-cN}$, where
$N=|V(G)|$, is $\#\mathrm{P}$-hard under deterministic polynomial-time
Turing reductions.  In particular, this applies to the tied-phase first
and second derivatives.
\end{corollary}

\begin{proof}%



The Fourier transform of the derivative oracle %
yields approximations to
\((ik)^r a_k\). As in Corollary~\ref{cor:precision}, for every \(k>D_\star\) the true
coefficient vanishes, whereas the coefficient at \(k=D_\star\) exceeds a
known computable threshold in modulus. Scanning the recovered
coefficients from the largest exponent downward therefore identifies
\(D_\star\).


Suppose that each derivative query is returned with additive error at most
$\varepsilon$.  The normalized inverse Fourier transform shows that the
recovered value of $(ik)^ra_k$ has additive error at most
$\varepsilon$.  Consequently, the recovered value of $a_{D_\star}(G)$ has
error at most
$\varepsilon/D_\star^r$. 
On the constructed family,
\[
  a_{D_\star}(G(H))
  =
  \kappa(H)\,|\M(H)|,
\]
with $|\kappa(H)|=2^{-O(N)}$ as shown. 
It is therefore enough to require
\[
  \frac{\varepsilon}{D_\star^r}
  <
  \frac{|\kappa(H)|}{2}.
\]
Because $D_\star=\operatorname{poly}(N)$ and $r$ is fixed, this inequality
holds for $\varepsilon\le2^{-cN}$ with a suitable constant~$c>0$.
Division by the known factor $\kappa(H)$ followed by rounding recovers the
exact count.
\end{proof}









\section{Discussion} \label{sec:discussion}

Our main theorem gives a clean exact-complexity classification by depth for the standard QAOA for MaxCut
expectation values on unrestricted graphs: $p=1$ is polynomial-time computable, whereas every
fixed $p\geq2$ is already $\#$P-hard.  Our reduction uses high-degree anchors, consistent
with the tractability of fixed-depth cost expectation value evaluation on bounded-degree or bounded-local-treewidth
graph families~\cite{WangEtAl2025}.  
\sh{This degree dependence differs sharply from
that of the full-distribution sampling problem: for general
$2$-local cost %
Hamiltonians, \={A}boli\c{n}\v{s} and
Ambainis~\cite{abolicnvs2026sharp} recently showed that depth-one QAOA
sampling is already classically hard 
for multiplicative-error sampling 
at interaction (maximum vertex) degree three,
whereas degree-two %
instances are exactly classically sampleable in
polynomial time for $p=O(\log n)$}. 
Thus the relevant degree
threshold depends essentially on whether the classical task is
evaluation of a local or global expectation value or sampling from
the complete QAOA output distribution. 
Our reduction also concerns exact or exponentially precise algebraic evaluation; no claims are
made about computations outputting constant or polynomially small error. 



Our main result is worst-case and %
does not imply hardness for all or for typical graphs (MaxCut problem instances). 
An interesting future direction is a more precise characterization of which graph subfamilies our hardness obstruction applies to or not. 
As mentioned, analytic or numerical solutions are possible in a number of cases. For example, a sequence of papers has obtained successively improved results concerning QAOA for MaxCut on high-girth graphs~\cite{Wang2018,hastings2019classical,wurtz2021maxcut,marwaha2021local,barak2021classical,basso2021quantum,sureshbabu2024parameter,farhi2025lower}, where our present reduction %
does not apply. 
\sh{A particularly important open direction is to determine %
to what degree computational hardness of QAOA expectation values persists
at physically relevant levels of precision, in particular for constant or inverse-polynomial
additive error, as well as under physical noise~\cite{marshall2020characterizing}. Resolving this question would help clarify the extent to which the
exact-complexity transition at $p=2$ constrains efficient quantum or classical %
evaluation 
of QAOA expectation values, both theoretically and in practice. Similar considerations apply to related quantities such as their gradients and Hessian important for parameter training.}


\paragraph{Implications for algorithms based on QAOA expectation values}
Recently a number of hybrid approaches have been proposed based on classical processing of quantum expectation values obtained using a quantum ansatz such as QAOA. 
Examples include so-called %
Iterative %
Quantum Algorithms%
~\cite{bravyi2020obstacles,bravyi2020hybrid,dupont2023quantum,MIS,Finzgar2023,brady2025quantum,wybo2026scalable}, where %
correlators $\langle Z_iZ_j\rangle$ %
are iteratively estimated and processed to obtain a sequence of reduced problem instances, %
and related relaxation-based approaches~\cite{dupont2024extending,maciejewski2024multilevel}, where the 
expectation values are classically processed to  produce candidate solutions directly. 
Strengthening the implications of~\cite{WangEtAl2025}, Theorem~\ref{thm:main} shows that the expectation values %
appearing in these and other related approaches can be computationally hard to evaluate in the worst case, at lease if exponential precision is required;  
we point the reader to \cite[Sec. III.D]{wybo2026scalable} for additional relevant discussion. %
Our present %
results neither establish nor rule out %
potential quantum
advantage for these algorithms. 








\paragraph{Connections to quantum supremacy experiments}
Our main result provides an %
expectation value analogue of existing sampling-hardness
results for low-depth QAOA  circuits~\cite{farhi2016quantum,harrow2017quantum,dalzell2020many}. %
Already at QAOA depth $p=2$, we show %
classical evaluation of a restricted QAOA for MaxCut expectation value %
suffices to recover
$\#\mathrm{P}$-hard counting information.
However, this does not by itself constitute an experimentally accessible
quantum-advantage result.  In particular, by Hoeffding's inequality estimating a bounded
observable to additive error $\varepsilon$ using independent circuit
measurements requires
$O(\varepsilon^{-2}\log\frac{1}{\delta})$
independent samples to achieve failure probability at most $\delta$.  
\sh{For the global MaxCut observable, a
computational basis measurement instead returns a cut value in
$[0,|E|]$.  Directly estimating the unnormalized
expectation $\langle C_G\rangle$ to additive error $\varepsilon$
therefore requires $O(
    m^2\varepsilon^{-2}\log\frac{1}{\delta}
  )$
samples.} 
The
accuracy $\varepsilon=2^{-\alpha N}$ used in our reduction %
hence entails exponentially many direct measurement samples; the additional
$m^2=\textrm{poly}(N)$ factor for the global objective does not change this
conclusion. 
Our result should consequently be interpreted as a structural
worst-case statement about %
QAOA
expectation value landscapes, rather than as %
one concerning 
experimentally accessible precision. %














\section*{Acknowledgment}
\sh{%
The author is grateful to Lucas Braydwood and Phillip Lotshaw for helpful discussions regarding the results of Ref.~\cite{WangEtAl2025}.} 
This material is based upon
work supported by the U.S. Department of Energy, Office of Science, Office of Advanced Scientific Computing Research under Award Number 89243024SSC000129, %
and under %
Prime Contract No. 80ARC020D0010 with the
NASA Ames Research Center. 
\sh{The author acknowledges the use of ChatGPT-5.6-Sol AI tools as part of the drafting of Figure~\ref{fig:triangle-gadget} and for verification of the proofs of this manuscript. %
The author assumes complete responsibility for all reported findings.} %


\bibliographystyle{unsrt}
\nocite{lenstra1982factoring,basu2006algorithms,cohen2013course} %added for footnote citation
\bibliography{bib}


\appendix







\section{Depth one QAOA expectation values are exactly tractable}\label{sec:p1}

For an edge $e=\{u,v\}$, let $d_u,d_v$ be the endpoint degrees and
$\lambda_{uv}=|N(u)\cap N(v)|$.  
For each single edge term $C_{uv}=(I-Z_uZ_v)/2$ 
the standard depth-one QAOA  calculation~\cite{Wang2018,hadfield2018quantum} gives
\begin{align}
 \langle C_{uv}\rangle_{p=1}
 ={}&\frac12
 +\frac14\sin(4\beta)\sin\gamma
 \left[(\cos\gamma)^{d_u-1}+(\cos\gamma)^{d_v-1}\right]
 \notag\\
 &-\frac14\sin^2(2\beta)
 (\cos\gamma)^{d_u+d_v-2\lambda_{uv}-2}
 \left[1-(\cos2\gamma)^{\lambda_{uv}}\right],
 \label{eq:p1-formula}
\end{align}
with  $F_1(G;\boldsymbol\gamma,\boldsymbol\beta)
 = \sum_{\{u,v\}\in E} \langle C_{uv}\rangle_{p=1}$.
This formula follows by conjugating $Z_uZ_v$ by the mixer and phase operators, and retaining only Pauli
strings with nonzero expectation values for the initial state $|+\rangle^{\otimes n}$. %

\begin{proposition}\label{prop:p1}
$\Eval^{uv}_1,\,\Eval^{\mathrm{MC}}_1\in\FP$.
\end{proposition}

\begin{proof}
Compute all degrees and common-neighbor counts, evaluate \eqref{eq:p1-formula} for every
edge, and sum, which requires a number of operations polynomially scaling in the problem size.  Exact exponentiation and arithmetic over the algebraic angle field require
polynomial bit complexity. 
\end{proof}

Our results above show that, under standard complexity-theoretic assumptions, no analogous exact efficiently computable formulas are possible for arbitrary graphs with $p\geq 2$ in general. 















\end{document}
